\documentclass{article}
\usepackage[T1]{fontenc} 
\usepackage[utf8]{inputenc}
\usepackage{graphicx}
\usepackage{xcolor}
\usepackage{float}
\usepackage[english]{babel}
\usepackage{rotating}
\usepackage[colorlinks]{hyperref}
\usepackage[english, boxed]{algorithm2e}
\usepackage{listings}
\definecolor{codegreen}{rgb}{0,0.6,0}
\definecolor{codegray}{rgb}{0.5,0.5,0.5}
\definecolor{codepurple}{rgb}{0.58,0,0.82}
\definecolor{backcolour}{rgb}{0.95,0.95,0.92}

\lstdefinestyle{mystyle}{
    backgroundcolor=\color{backcolour},   
    commentstyle=\color{codegreen},
    keywordstyle=\color{magenta},
    numberstyle=\tiny\color{codegray},
    stringstyle=\color{codepurple},
    basicstyle=\ttfamily\footnotesize,
    breakatwhitespace=false,         
    breaklines=true,                 
    captionpos=b,                    
    keepspaces=true,                 
    numbers=left,                    
    numbersep=5pt,                  
    showspaces=false,                
    showstringspaces=false,
    showtabs=false,                  
    tabsize=2
}

\lstset{style=mystyle}

\begin{document}
    %%%%%%%%%%%%%%%%%%%%%%%%% TITLE PAGE %%%%%%%%%%%%%%%%%%%%%%%%%
    \begin{titlepage}
        \centering
        \includegraphics[scale=0.1,keepaspectratio]{img/logo_en.png}
        \vspace{5cm}
        
        {\textbf{\Huge Security system based on ESP32}}
        
        \vspace{0.25cm}
        {\LARGE 2021/2022}
        \vspace{1cm}
        
        {\textbf{\LARGE Microprocessors and Embedded Systems}}
        \vspace{0.25cm}
        
        
        {\LARGE Project documentation}
        \vspace{1cm}
        
        {\Large Roman Janiczek (xjanic25)}
        
        \vspace{1cm}
        {\large \today}
    
    \end{titlepage}
    
    \newpage
    
    %%%%%%%%%%%%%%%%%%%%%%%%%%% CONTENT %%%%%%%%%%%%%%%%%%%%%%%%%%
    \clearpage
    \tableofcontents
    \thispagestyle{empty}
    %%%%%%%%%%%%%%%%%%%%%%%% INTRODUCTION %%%%%%%%%%%%%%%%%%%%%%%%
    \clearpage
    \pagenumbering{arabic}
    \setcounter{page}{1}
    
    \section{Introduction}
    This project aims to create a simple security system that 
    consists of one PIR sensor and a camera. The system should also
    have WiFi AP capabilities to view captured pictures on the 
    webserver running in this internal network.
    
    For this, we will be using \verb|WeMos D1 R32 UNO ESP32|,    \\
    \verb|AM312 PIR Motion sensor| and \verb|AI-Thinker's ESP32-CAM|.
    
    \subsection{WeMos D1 R32 UNO ESP32}
    \verb|WeMos D1 R32 UNO ESP32| is par of ESP32 series.
    This is a series low-cost, low-power MCUs with integrated Wi-Fi
    and Bluetooth.\footnote{\url{https://en.wikipedia.org/wiki/ESP32}}
    
    \subsection{AI-Thinker's ESP32-CAM}
    \verb|AI-Thinker's ESP32-CAM| is ESP32-based small sized camera 
    with integrated Wi-Fi and support for TF cards.\footnote{\url{https://docs.ai-thinker.com/en/esp32-cam}}
    
    %%%%%%%%%%%%%%%%%%%%%%%%%%% DESIGN %%%%%%%%%%%%%%%%%%%%%%%%%%%
    \section{Design}
    The system will consist of two parts that will be interconnected.
    \begin{itemize}
        \item Motion detection
        \item Camera and AP
    \end{itemize}
    \begin{figure}[h!]
        \centering
        \includegraphics[scale=0.43,keepaspectratio]{img/circuit.png}
        \caption{Connection diagram}\label{pic:connection_diagram}
    \end{figure} 
    
    \subsection{Motion detection}
    Motion detection will be realized using \verb|AM312 PIR Motion sensor|\\
    and a \verb|WeMos D1 R32 UNO ESP32|.
    
    The motion detection model attempts to connect to
    a local WiFi created by a camera module as a first step. 
    This will allow it to send information about a detected motion 
    to the camera.
    
    Detected motion sends a signal from a PIR sensor to an ESP32 unit
    that creates an HTTP request for a second ESP32 unit (Camera), 
    informing it about this event.
    
    \subsubsection{Process of detecting movement}
    \begin{itemize}
        \item Wait for a movement (loop waiting)
        \item Receive signal from a PIR sensor
        \item Blink debug LED to let us know about detection
        \item Send signal over HTTP to camera
        \item Process response and go back to waiting
    \end{itemize}
    
    \subsection{Camera and AP}
    The camera module is responsible for creating a local WiFi network
    . This module also serves as a web server that
    displays the latest captured photo.
    
    Another functionality of the camera module and its primary purpose is
    capturing and saving a photo after receiving the signal 
    from a motion detection module.
    
    \subsubsection{Process of capturing a photo}
    \begin{itemize}
        \item Receive signal over HTTP about a movement
        \item Capture a photo (do not retry if failed)
        \item Save a photo to a memory
        \item Send a response to the motion detection module
            (information about the status of photo capture)
    \end{itemize}
    
    \subsubsection{Process of presenting a photo}
    \begin{itemize}
        \item Listens for requests to a web server
        \item Receive a request for a latest photo
        \item Load the photo from a memory
        \item Display it on the main page
            (or send to a requester depends on request
            URL)
    \end{itemize}
    
    %%%%%%%%%%%%%%%%%%%%%%% IMPLEMENTATION %%%%%%%%%%%%%%%%%%%%%%%
    \section{Implementation}
    
    \subsection{Used technologies}
    Implementation of this system was realized using 
    \href{https://docs.espressif.com/projects/esp-idf/en/latest/esp32/}{ESP-IDF}
    \footnote{\url{https://docs.espressif.com/projects/esp-idf/en/latest/esp32/}}
    , programming language \href{https://www.iso.org/standard/74528.html}{C}
    \footnote{\url{https://www.iso.org/standard/74528.html}} and 
    \href{https://docs.platformio.org/en/latest/integration/ide/vscode.html}{PlatformIO in VSCode}
    \footnote{\url{https://docs.platformio.org/en/latest/integration/ide/vscode.html}}.
    
    \subsection{Motion detection}
    \subsubsection{PIR sensor}
    GPIO first had to be chosen and configured as an input to use PIR sensor.
    For this, GPIO39 was selected. At the same time, debug LED output was configured
    on GPIO2.
    \begin{lstlisting}[language=C]
#include <driver/gpio.h>

#define LED_GPIO GPIO_NUM_2
#define PIR_GPIO GPIO_NUM_39

// Setup OUTPUT
gpio_pad_select_gpio(LED_GPIO);
gpio_set_direction(LED_GPIO, GPIO_MODE_OUTPUT);
gpio_set_level(LED_GPIO, 0);
// Setup INPUT
gpio_pad_select_gpio(PIR_GPIO);
gpio_set_direction(PIR_GPIO, GPIO_MODE_INPUT);
    \end{lstlisting}
    
    After correct configuration of INPUT GPIO it is left running
    in infinite while loop where it checks if a signal has been received.
    If it detects signal it blinks with debug LED and calls function
    to send HTTP request to camera.
    \begin{lstlisting}[language=C]
while(1) {  
    if (gpio_get_level(PIR_GPIO)) {
        timestamp = esp_timer_get_time();
        gpio_set_level(LED_GPIO, 1);
        send_request_to_camera();
        vTaskDelay(delay);
        gpio_set_level(LED_GPIO, 0);
    } else {
        gpio_set_level(LED_GPIO, 0);
        vTaskDelay(delay);
    }
}
    \end{lstlisting}
    
    To prevent it from sending too much requests processing is delayed
    for 1s using function \verb|vTaskDelay()| and const delay.
    \begin{lstlisting}[language=C]
const TickType_t delay = 1000 / portTICK_PERIOD_MS;
    \end{lstlisting}
    
    
    \subsubsection{Wireless connection (WiFi STA)}
    Connection to existing WiFi AP is realized using library \verb|"esp_wifi.h"|
    and official example code\footnote{\url{https://github.com/espressif/esp-idf/blob/5c33570524118873f7bd32490c7a0442fede4bf8/examples/wifi/fast_scan/main/fast_scan.c}}.
    Default SSID for AP is hardcoded. Our IP and a gateway are
    saved for further use later in a program during connection.
    \begin{lstlisting}[language=C]
#define WIFI_SSID       "ESP32-Cam AP"
.
my_ip = event->ip_info.ip;
gateway = event->ip_info.gw;
    \end{lstlisting}
    
    \subsubsection{HTTP requests}
    Sending of a signal to a camera is realized using HTTP request. This functionality
    can be implemented using library \verb|"esp_http_client.h"|. For this project was
    used modified official example code\footnote{\url{https://github.com/espressif/esp-idf/blob/5c33570524118873f7bd32490c7a0442fede4bf8/examples/protocols/esp_http_client/main/esp_http_client_example.c}}.
    Here comes a previously saved gateway address into play - in this system, the gateway is always
    the camera module. The gateway address is taken as a host and appended \verb|/pir| path.
    HTTP client is configured with this newly created URL, initiated and called. 
    (Called only on PIR event)
    \begin{lstlisting}[language=C]
char temp_url[50];
sprintf(temp_url, "http://%d.%d.%d.%d/take-photo", IP2STR(&gateway));
    
esp_http_client_config_t config = {
    .url = temp_url,
    .method = HTTP_METHOD_GET,
    \end{lstlisting}

    \begin{lstlisting}[language=C]
esp_http_client_handle_t client = esp_http_client_init(&config);
esp_err_t err = esp_http_client_perform(client);
    \end{lstlisting}
    
    An HTTP request uses the method GET due to problems encountered with the POST method.
    
    \clearpage
    \subsection{Camera and AP}
    \subsubsection{Wireless connection (WiFi AP)}
    Camera module servers as the main hub of the system as such it's this module
    that creates WiFi AP. Library \verb|"esp_wifi.h"| is used here as well but this time
    during configuration we chose AP mode.
    \begin{lstlisting}[language=C]
esp_netif_create_default_wifi_ap();
    \end{lstlisting}
    \begin{lstlisting}[language=C]
wifi_config_t custom_config = {
    .ap = {
        .ssid = WIFI_SSID,
        .ssid_len = strlen(WIFI_SSID),
        .channel = WIFI_CHAN,
        .max_connection = WIFI_MSCO,
        .authmode = WIFI_AUTH_OPEN      // Do not require password
    },
};
ESP_ERROR_CHECK(esp_wifi_set_mode(WIFI_MODE_AP));
ESP_ERROR_CHECK(esp_wifi_set_config(WIFI_IF_AP, &custom_config));
ESP_ERROR_CHECK(esp_wifi_start());
    \end{lstlisting}
    
    \subsubsection{Storage}
    Due to problems encountered when trying to store pictures on an SD card or SPIFFS
    \footnote{\url{https://docs.espressif.com/projects/esp-idf/en/latest/esp32/api-reference/storage/spiffs.html}},
     I decided to implement the last taken photo storage as a temporary global
     variable that is rewritten when a new photo is taken. This decision also makes it easier
     to serve this photo on the webserver.
     \begin{lstlisting}[language=C]
/**
 * Latest photo store 
 */
size_t taken_photo_len = 0;
size_t taken_photo_width = 0;
size_t taken_photo_height = 0;
pixformat_t taken_photo_format = PIXFORMAT_JPEG;
uint8_t *taken_photo_buf = NULL; 
    \end{lstlisting}
    \begin{lstlisting}[language=C]
if (taken_photo_buf != NULL) {
    free(taken_photo_buf);
}
taken_photo_buf = malloc(photo->len);
mempcpy(taken_photo_buf, photo->buf, photo->len);

taken_photo_format = photo->format;
taken_photo_len = photo->len;
taken_photo_width = photo->width;
taken_photo_height = photo->height;
    \end{lstlisting}
    
    \subsubsection{Camera}
    To initialize the camera, a correct pin schema was needed. This was obtained from ESP-IDF GitHub\footnote{\url{https://github.com/espressif/esp32-camera/blob/master/examples/main/take_picture.c}}.
    I have also decided to enable flash when camera is taking pictures which allows for better visibility
    during night or in dark rooms.
    \begin{lstlisting}[language=C]
gpio_set_level(4, 1);
vTaskDelay(delay);
camera_fb_t *photo = esp_camera_fb_get();
gpio_set_level(4, 0);
    \end{lstlisting}
    
    \subsubsection{Webserver}
    Webserver as de facto main 'brain' of the whole system which allows the processing of PIR event and is also
    serving latest taken photo is implemented using \verb|"esp_http_server.h"|\footnote{\url{https://docs.espressif.com/projects/esp-idf/en/latest/esp32/api-reference/protocols/esp_http_server.html}} library from ESP-IDF.
    
    The webserver is served on the default IP address of the camera module \verb|192.168.4.1| and allows for the following paths:
    \begin{itemize}
        \item \verb|/| - main page which displays latest taken photo and automatically refreshes every 5 seconds
        \item \verb|/latest-photo.jpg| - server latest photo captured by PIR event directly
        \item \verb|/take-photo| - captures a new photo and displays it directly; this photo is not saved anywhere and does not replace
        the latest photo captured by the PIR event
        \item \verb|/pir| - is used by motion sensor module as an endpoint for sending a signal about the PIR event
    \end{itemize}
    
    All of those endpoints are calling handler functions that are responsible for providing correct data.
    \begin{lstlisting}[language=C]
httpd_uri_t index_get = {
    .uri      = "/",
    .method   = HTTP_GET,
    .handler  = get_handler,
    .user_ctx = NULL
};
httpd_register_uri_handler(server, &index_get);
httpd_uri_t latest_get = {
    .uri      = "/latest-photo.jpg",
    .method   = HTTP_GET,
    .handler  = img_handler,
    .user_ctx = NULL
};
httpd_register_uri_handler(server, &latest_get);
httpd_uri_t pir_post = {
    .uri      = "/pir",
    .method   = HTTP_GET,
    .handler  = pir_handler,
    .user_ctx = NULL
};
httpd_register_uri_handler(server, &pir_post);
httpd_uri_t take_post = {
    .uri      = "/take-photo",
    .method   = HTTP_GET,
    .handler  = take_handler,
    .user_ctx = NULL
};
httpd_register_uri_handler(server, &take_post);
    \end{lstlisting}
    
    %%%%%%%%%%%%%%%%%%%%%%%%% CONCLUSION %%%%%%%%%%%%%%%%%%%%%%%%%
    \clearpage
    \section{Conclusion}
    In conclusion this project was a lot of fun as it allowed me to learn a lot about 
    ESP32 lower lever programming as compared to Arduino framework.
    The only thing that should be fixed that I'm aware of is the lack of archiving older photos on SD card or some other media.
    
    \subsubsection{Notes}
    Due to problem with programmer for \verb|AI-Thinker's ESP32-CAM| which cause following error to appear a lot even 
    when following correct procedure, programming and testing anything on \verb|AI-Thinker's ESP32-CAM| was incredibly difficult and time consuming.
    \begin{lstlisting}
Traceback (most recent call last):
File "C:\Users\roman\.platformio\packages\tool-esptoolpy\esptool.py", line 4582, in <module>
    _main()
File "C:\Users\roman\.platformio\packages\tool-esptoolpy\esptool.py", line 4575, in _main
    main()
File "C:\Users\roman\.platformio\packages\tool-esptoolpy\esptool.py", line 4074, in main
    esp = esp or get_default_connected_device(ser_list, port=args.port, connect_attempts=args.connect_attempts,
File "C:\Users\roman\.platformio\packages\tool-esptoolpy\esptool.py", line 121, in get_default_connected_device
    _esp.connect(before, connect_attempts)
File "C:\Users\roman\.platformio\packages\tool-esptoolpy\esptool.py", line 632, in connect
    last_error = self._connect_attempt(mode=mode, esp32r0_delay=True, usb_jtag_serial=usb_jtag_serial)
File "C:\Users\roman\.platformio\packages\tool-esptoolpy\esptool.py", line 596, in _connect_attempt
    self.sync()
File "C:\Users\roman\.platformio\packages\tool-esptoolpy\esptool.py", line 486, in sync
    val, _ = self.command(self.ESP_SYNC, b'\x07\x07\x12\x20' + 32 * b'\x55',
File "C:\Users\roman\.platformio\packages\tool-esptoolpy\esptool.py", line 418, in command
    self._port.timeout = new_timeout
File "C:\Users\roman\.platformio\penv\lib\site-packages\serial\serialutil.py", line 372, in timeout
    self._reconfigure_port()
File "C:\Users\roman\.platformio\penv\lib\site-packages\serial\serialwin32.py", line 222, in _reconfigure_port
    raise SerialException(
serial.serialutil.SerialException: Cannot configure port, something went wrong. Original message: PermissionError(13, 'Access is denied.', None, 5)
*** [upload] Error 1
    \end{lstlisting}
    
    %%%%%%%%%%%%%%%%%%%%%%%% ATTACHMENTS %%%%%%%%%%%%%%%%%%%%%%%%%
    \clearpage
    \section{Attachments}
    \begin{figure}[h!]
        \centering
        \includegraphics[scale=0.35,keepaspectratio]{img/webservermain.png}
        \caption{Main page of webserver}\label{pic:webserverindex}
    \end{figure} 
    \begin{figure}[h!]
        \centering
        \includegraphics[scale=0.43,keepaspectratio]{img/pirevent.png}
        \caption{Logged PIR event on motion detection module}\label{pic:pireventlog}
    \end{figure} 
    %%%%%%%%%%%%%%%%%%%%%%%% BIBLIOGRAPHY %%%%%%%%%%%%%%%%%%%%%%%%
    \clearpage
    %bibliographystyle{englishiso}
    %renewcommand{\refname}{Bibliography}
    %bibliography{documentation}
	
\end{document}